\newpage
\section{Testare și rezultate}

Acest capitol prezintă metodologia de testare aplicată sistemului și rezultatele obținute. Sunt urmărite patru aspecte: corectitudinea funcțională, performanța modelului de inteligență artificială, latența sistemului și fiabilitatea mecanismului de alertare.

\textit{[Notă: valorile numerice marcate mai jos cu paranteze drepte trebuie completate cu rezultatele măsurate efectiv în timpul testării. Metodologia este descrisă complet; doar cifrele finale rămân de inserat.]}

\subsection{Metodologia de testare}

Testarea sistemului a fost realizată pe mai multe planuri. \textbf{Testarea funcțională} a vizat verificarea fiecărui flux descris în subcapitolul 4.4, anume înregistrarea utilizatorilor, asocierea unei persoane monitorizate, recepția măsurătorilor, declanșarea alertelor și livrarea notificărilor.

Pentru testarea fluxurilor care depind de dispozitivul fizic a fost utilizat un \textbf{mod de simulare}, integrat în aplicație, care permite generarea de măsurători cu valori controlate, fără a fi necesar hardware-ul real. Acest mecanism a permis reproducerea sistematică a scenariilor de risc, anume valori critice de \spo, tahicardie, febră și imobilitate, urmate de verificarea reacției corecte a sistemului în fiecare caz.

\textbf{Performanța modelului de inteligență artificială} a fost evaluată separat, pe setul de date, folosind tehnicile de validare descrise mai jos.

\subsection{Testare automată}

Pe lângă testarea funcțională manuală, aplicația dispune de o suită de \textbf{180 de teste automate}, organizate în două categorii, executate automat la fiecare push în pipeline-ul CI/CD configurat cu GitHub Actions.

\subsubsection{Teste unitare}

Testele unitare verifică izolat comportamentul fiecărui serviciu și controler, folosind obiecte mock (prin biblioteca \textbf{Moq}) pentru a elimina dependențele externe (baza de date, email, SMS). Sunt acoperite: \texttt{AuthenticationService} (înregistrare, autentificare, verificare email, resetare parolă), \texttt{JwtService} (generare și validare token), \texttt{UserService}, \texttt{MeasurementService}, \texttt{MonitoredService}, \texttt{WifiNetworkService} și \texttt{ConditionRuleEngine} (motorul de reguli medicale deterministe). La nivel de controlere REST sunt testate scenariile principale de succes și de eșec pentru \texttt{AuthenticationController}, \texttt{MeasurementController}, \texttt{WifiController} și \texttt{ConfigController}.

Aserțiunile sunt formulate cu biblioteca \textbf{FluentAssertions}, care produce mesaje de eroare explicite în caz de eșec, facilitând identificarea cauzei fără a interpreta stack trace-uri brute.

\subsubsection{Teste de integrare}

Testele de integrare verifică interacțiunea dintre servicii și baza de date, folosind o instanță \textbf{SQLite in-memory} izolată per test (fiecare test pornește cu o bază curată, creată prin \texttt{EnsureCreated}). Sunt acoperite: \texttt{UserRepository} (CRUD utilizatori, autentificare, confirmare email), \texttt{MeasurementRepository} (stocare și interogare măsurători cu filtre temporale), \texttt{WifiNetworkRepository} (adăugare și ștergere rețele WiFi per dispozitiv), \texttt{RetentionCleanupService} (verificarea că job-ul de curățare elimină corect datele mai vechi decât perioada de retenție) și \texttt{NotificationFeedbackService} (ciclul complet de creare, marcare ca \textit{așteptând feedback} și înregistrare a răspunsului utilizatorului).

\subsubsection{Rezultate}

Toate cele 180 de teste trec cu succes în configurația curentă (0 eșuate, 0 sărite). Durata totală de execuție a suitei complete este de aproximativ 24 de secunde, acceptabilă pentru un ciclu CI/CD. Integrarea în GitHub Actions asigură că orice regresie este detectată înainte de publicarea codului în producție.

\subsection{Acuratețea datelor și performanța modelului AI}

Performanța modelului Random Forest a fost estimată prin \textbf{validare încrucișată stratificată în 5 partiții}, metodă care împarte repetat setul de date în subseturi de antrenare și testare, păstrând în fiecare partiție distribuția claselor NORMAL, ALERT și CRITICAL. Au fost urmărite acuratețea, scorul F1 (în variantele macro și ponderată) și aria de sub curba ROC (AUC-ROC), calculată în regim \textit{one-vs-rest}.

Setul de date utilizat la antrenare și evaluare a fost construit sintetic, pornind de la sursele publice menționate în subcapitolul 2.3. El cuprinde câteva mii de exemple corespunzătoare stării normale și câteva mii de exemple corespunzătoare stărilor ALERT și CRITICAL, generate în mod echilibrat astfel încât modelul să dispună de suficiente cazuri din fiecare clasă. Caracterul stratificat al validării garantează că această proporție este păstrată în fiecare partiție, evitând estimările optimiste care ar putea apărea dacă o clasă rară ar fi insuficient reprezentată în subsetul de testare.

Rezultatele obținute sunt sintetizate în tabelul~\ref{tab:cv-results}.

\begin{table}[H]
\centering
\small
\renewcommand{\arraystretch}{1.3}
\begin{tabular}{|>{\raggedright\arraybackslash}p{5.5cm}|c|c|}
\hline
\textbf{Metrică} & \textbf{Medie} & \textbf{Abatere standard}\\
\hline
Acuratețe & [\ \ ] & [\ \ ]\\
\hline
Scor F1 (macro) & [\ \ ] & [\ \ ]\\
\hline
Scor F1 (ponderat) & [\ \ ] & [\ \ ]\\
\hline
AUC-ROC (one-vs-rest) & [\ \ ] & N/A\\
\hline
\end{tabular}
\caption{Rezultatele validării încrucișate stratificate în 5 partiții.}
\label{tab:cv-results}
\end{table}

Pentru a verifica absența supraînvățării (overfitting), a fost comparată acuratețea pe datele de antrenare cu cea pe datele de testare; o diferență mică între cele două indică o bună capacitate de generalizare a modelului \textit{(diferența măsurată: [\ \ ])}.

Comportamentul modelului pe fiecare clasă a fost analizat prin \textbf{matricea de confuzie} (figura~\ref{fig:confuzie}), care evidențiază eventualele confuzii între stările apropiate. Importanța parametrilor de intrare a fost evaluată prin metoda \textbf{SHAP} (figura~\ref{fig:shap}), care a confirmat că saturația oxigenului, frecvența cardiacă și temperatura corporală sunt factorii dominanți în decizia modelului, în concordanță cu criteriile clinice de triaj.

\begin{figure}[H]
\centering
\figplaceholder{[Spațiu pentru figură: matricea de confuzie a modelului pe cele trei clase (NORMAL, ALERT, CRITICAL), generată de scriptul de antrenare.]}
\caption{Matricea de confuzie a modelului Random Forest.}
\label{fig:confuzie}
\end{figure}

\begin{figure}[H]
\centering
\figplaceholder{[Spațiu pentru figură: graficul SHAP cu importanța globală a parametrilor de intrare, generat de scriptul de antrenare.]}
\caption{Importanța parametrilor de intrare (analiză SHAP).}
\label{fig:shap}
\end{figure}

\subsection{Latența sistemului}

Latența reprezintă intervalul de timp scurs între achiziția unei măsurători pe dispozitiv și afișarea sau alertarea acesteia (obiectivul O5). Au fost măsurate două intervale: timpul \textit{dispozitiv $\rightarrow$ afișare}, de la citirea senzorilor până la actualizarea tabloului de bord, și timpul \textit{dispozitiv $\rightarrow$ alertă}, de la achiziție până la livrarea notificării push.

Latența totală este suma mai multor componente: intervalul de transmitere configurat pe dispozitiv, timpul de procesare pe backend, timpul de răspuns al microserviciului AI și timpul de livrare prin SignalR. Valorile măsurate au fost: timp \textit{dispozitiv $\rightarrow$ afișare} de aproximativ [\ \ ] secunde și timp \textit{dispozitiv $\rightarrow$ alertă} de aproximativ [\ \ ] secunde, confirmând funcționarea în regim apropiat de timp real.

\subsection{Fiabilitatea alertelor}

Fiabilitatea mecanismului de alertare a fost evaluată din două perspective: capacitatea de a semnala situațiile reale de risc și capacitatea de a evita alertele false. Folosind modul de simulare, au fost generate scenarii de risc cunoscute și s-a verificat dacă sistemul declanșează alerta corespunzătoare. În paralel, au fost generate valori normale cu variații tranzitorii, pentru a verifica dacă pragul de persistență elimină alertele false.

Rezultatele au arătat că, din [\ \ ] scenarii de risc testate, sistemul a declanșat corect alerta în [\ \ ] cazuri, iar pragul de persistență a redus numărul de alerte false cu [\ \ ]\% față de o variantă fără acest mecanism. A fost testat separat și comportamentul la pierderea conexiunii: măsurătorile acumulate în coada locală a dispozitivului au fost transmise integral și în ordine corectă la revenirea semnalului, fără pierderi de date.
